% =========================================================
% Density of the Rational Numbers
% =========================================================

\subsection{Density of the Rational Numbers}
\label{sec:density-q}

\subsection{Definitions and Theorems}

Two different ideas are often confused when one first studies the rational numbers.

\begin{itemize}
  \item The \emph{Archimedean property} says that the natural numbers grow beyond every fixed
  number.
  \item The \emph{density of $\mathbb{Q}$} says that between any two distinct real numbers there is a
  rational number.
\end{itemize}

The first controls large-scale size; the second controls local refinement of intervals.

\begin{theorem}[Archimedean property of $\mathbb{Q}$]
\label{thm:archimedean-q}
For every $x\in\mathbb{Q}$ there exists $n\in\mathbb{N}$ such that
\[
n>x.
\]
Equivalently, for every $x>0$ there exists $n\in\mathbb{N}$ such that
\[
\frac{1}{n}<x.
\]
\end{theorem}

\begin{remark*}[Proof]
\hyperref[prf:archimedean-q]{Go to proof.}
\end{remark*}

\begin{remark}
This theorem says that the reciprocals of natural numbers become arbitrarily small.  It rules out the
existence of positive infinitesimals in $\mathbb{Q}$.
\end{remark}

\begin{remark}[Forward reference]
The statement that $\mathbb{Q}$ is dense in $\mathbb{R}$ is best treated as a
theorem about how the rationals sit inside the completed real line, not as an
internal theorem about $\mathbb{Q}$ alone.  Its canonical home is therefore in
the reals chapter:
\hyperref[thm:q-dense-in-r]{Density of $\mathbb{Q}$ in $\mathbb{R}$}.  The
same is true of the corollary that there are infinitely many rationals between
two distinct real numbers.
\end{remark}

\subsection{Consequences}

\begin{remark}
Density is one of the reasons the rational numbers are so useful in analysis.  Even when a number is
not rational, it can still be approximated arbitrarily well by rational numbers.
\end{remark}
